\section{Forces}

A force is a rule for how a tone must rotate as it moves from dock to dock. The allowed rotations come from the number
systems nested inside the octonion, and there are exactly three nontrivial ones, giving exactly three forces.

The complex numbers give the symmetry of the unit circle, a one-parameter rotation, which is the photon and the
hypercharge. The quaternions give the symmetry of the unit three-sphere, a three-parameter rotation, which is the weak
force. And the octonions, through their automorphism group, give the eight-parameter color symmetry as the part that
fixes one imaginary unit, which is the strong force. Together this is the gauge group of the Standard Model, read off the
three division algebras of the tower, rung by rung. The explicit generators and the dimension counts are computed and
checked in the open codebase \cite{cluesurf2026repo}.

The three forces behave as they should. The photon is massless because its symmetry is never broken, so it reaches across
the whole mesh undimmed. The weak force is short-range because the vacuum's chosen alignment, the Higgs, clings to it
and gives it mass. The strong force traps itself, because its eight carriers themselves carry color, so they rotate one
another and bunch up rather than spreading, which is why a free quark is never seen. That confinement is reproduced on
the substrate as a real string tension at every coupling, two opposite charges in one dimension feeling a constant force
that never lets them part.

So the three forces of nature are the three nontrivial division algebras seen as symmetries, and the matter they act on,
from the previous section, is the spinor built from the same octonion. One algebra gives both the players and the rules.
The exact nature of the gauge group, and the one place it admits two readings, is the next section.
